% =============================================================================
%  3.1 Cartografía
% =============================================================================

\subsection{Visualización Cartográfica de Resultados}

% ---
\begin{frame}{Isócronas peatonales}
  \begingroup
    \centering
    \includegraphics[width=\linewidth,height=0.63\textheight,keepaspectratio]
    {Figures/Qgis/Mapa_Cobertura_Red_Total.pdf}\par
  \endgroup

  \footnotesize
  \begin{itemize}
    \item Delimitan la Cobertura Espacial ($C$) real sobre la red vial, exponiendo 
          la drástica caída de acceso al salir de la CDMX.
  \end{itemize}
\end{frame}

% ---
\begin{frame}{Fuerza Capilar ($FC$)}
  \begingroup
    \centering
    \includegraphics[width=\linewidth,height=0.63\textheight,keepaspectratio]
    {Figures/Qgis/Mapa_Fuerza_Capilar.pdf}\par
  \endgroup

  \footnotesize
  \begin{itemize}
        \item Muestran el grado nodal de Fuerza Capilar ($FC$) y los puntos de 
        Centralidad de Intermediación ($B(v)$), coloreados según las cuatro bandas de la \termino{Escala Garibelt}.
  \end{itemize}
\end{frame}
